% GENERATED FILE. Edit canonical lessons, then run npm run generate:books.
% canonical-source-hash: fnv1a64:846755478a1234c8
% canonical-lessons: KA-C118-kannadaka, KA-C118-gadiyara, KA-C118-cila, KA-C118-lekhani, KA-C118-kagada

\chapter{Glasses, Clock, Bag, Pen, Paper}
\label{ch:ka-glasses-clock-bag-pen-paper}
\hlchaptermodality{118}

\begin{chapteropening}
\textbf{By the end of this chapter:} \emph{I can say glasses, a clock, a bag, a pen and paper.}

Complete the last lesson of chapter 118: i can say glasses, a clock, a bag, a pen and paper.
\end{chapteropening}

\section[\texorpdfstring{\kn{ಕನ್ನಡಕ} (kannaḍaka)}{ಕನ್ನಡಕ (kannaḍaka)}]{\kn{ಕನ್ನಡಕ} (kannaḍaka) — glasses, spectacles}

\label{lesson:KA-C118-kannadaka}



\begin{quote}
Before the new one: say the Kannada for soap, then the Kannada for a towel.
\end{quote}

\subsection*{You'll want to know: \kn{ಕನ್ನಡಕ}}
\textbf{\kn{ಕನ್ನಡಕ}} — \emph{kannaḍaka} — \textquotedblleft{}glasses, spectacles\textquotedblright{}.

\begin{grammarlens}[title={What you've built}]
One more word for this chapter.
\end{grammarlens}

\subsection*{Your turn}
\begin{itemize}
\raggedright
  \item \emph{Say it:} \emph{kannaḍaka}
  \item \emph{Say it:} \emph{kannaḍaka}, once more
  \item \emph{Say it:} say \emph{kannaḍaka}
  \item \emph{Recall:} say the Kannada for soap, then the Kannada for a towel, then say \emph{kannaḍaka} again
\end{itemize}

\begin{tcolorbox}[breakable,colback=teal!4,colframe=teal!35!black,title={Before you move on}]
What does \kn{ಕನ್ನಡಕ} mean? (\textquotedblleft{}Glasses, spectacles\textquotedblright{}.) Say it once more.
\end{tcolorbox}
\section[\texorpdfstring{\kn{ಗಡಿಯಾರ} (gaḍiyāra)}{ಗಡಿಯಾರ (gaḍiyāra)}]{\kn{ಗಡಿಯಾರ} (gaḍiyāra) — a clock, a watch}

\label{lesson:KA-C118-gadiyara}



\begin{quote}
Before the new one: say the Kannada for a towel, then the Kannada for glasses.

Read the digits aloud: \textbf{\kn{೨}}, \textbf{\kn{೪}}, \textbf{\kn{೫}}, \textbf{\kn{೬}}, \textbf{\kn{೭}}, \textbf{\kn{೮}}, \textbf{\kn{೯}}, \textbf{\kn{೦}}.
\end{quote}

\subsection*{You'll want to know: \kn{ಗಡಿಯಾರ}}
\textbf{\kn{ಗಡಿಯಾರ}} — \emph{gaḍiyāra} — \textquotedblleft{}a clock, a watch\textquotedblright{}.

\begin{grammarlens}[title={What you've built}]
One more word for this chapter.
\end{grammarlens}

\subsection*{Your turn}
\begin{itemize}
\raggedright
  \item \emph{Say it:} \emph{gaḍiyāra}
  \item \emph{Say it:} \emph{gaḍiyāra}, once more
  \item \emph{Say it:} say \emph{gaḍiyāra}
  \item \emph{Recall:} say the Kannada for a towel, then the Kannada for glasses, then say \emph{gaḍiyāra} again
\end{itemize}

\begin{tcolorbox}[breakable,colback=teal!4,colframe=teal!35!black,title={Before you move on}]
What does \kn{ಗಡಿಯಾರ} mean? (\textquotedblleft{}A clock, a watch\textquotedblright{}.) Say it once more.
\end{tcolorbox}
\section[\texorpdfstring{\kn{ಚೀಲ} (cīla)}{ಚೀಲ (cīla)}]{\kn{ಚೀಲ} (cīla) — a bag}

\label{lesson:KA-C118-cila}



\begin{quote}
Before the new one: say the Kannada for glasses, then the Kannada for a clock.
\end{quote}

\subsection*{You'll want to know: \kn{ಚೀಲ}}
\textbf{\kn{ಚೀಲ}} — \emph{cīla} — \textquotedblleft{}a bag\textquotedblright{}.

\begin{grammarlens}[title={What you've built}]
One more word for this chapter.
\end{grammarlens}

\subsection*{Your turn}
\begin{itemize}
\raggedright
  \item \emph{Say it:} \emph{cīla}
  \item \emph{Say it:} \emph{cīla}, once more
  \item \emph{Say it:} say \emph{cīla}
  \item \emph{Recall:} say the Kannada for glasses, then the Kannada for a clock, then say \emph{cīla} again
\end{itemize}

\begin{tcolorbox}[breakable,colback=teal!4,colframe=teal!35!black,title={Before you move on}]
What does \kn{ಚೀಲ} mean? (\textquotedblleft{}A bag\textquotedblright{}.) Say it once more.
\end{tcolorbox}
\section[\texorpdfstring{\kn{ಲೇಖನಿ} (lēkhani)}{ಲೇಖನಿ (lēkhani)}]{\kn{ಲೇಖನಿ} (lēkhani) — a pen}

\label{lesson:KA-C118-lekhani}



\begin{quote}
Before the new one: say the Kannada for a clock, then the Kannada for a bag.

Say what the colon \textbf{:} and the brackets \textbf{( )} tell a reader. Then find \textbf{\kn{ಠ}} and say \emph{ṭha}.
\end{quote}

\subsection*{You'll want to know: \kn{ಲೇಖನಿ}}
\textbf{\kn{ಲೇಖನಿ}} — \emph{lēkhani} — \textquotedblleft{}a pen\textquotedblright{}.

\begin{grammarlens}[title={What you've built}]
One more word for this chapter.
\end{grammarlens}

\subsection*{Your turn}
\begin{itemize}
\raggedright
  \item \emph{Say it:} \emph{lēkhani}
  \item \emph{Say it:} \emph{lēkhani}, once more
  \item \emph{Say it:} say \emph{lēkhani}
  \item \emph{Recall:} say the Kannada for a clock, then the Kannada for a bag, then say \emph{lēkhani} again
\end{itemize}

\begin{tcolorbox}[breakable,colback=teal!4,colframe=teal!35!black,title={Before you move on}]
What does \kn{ಲೇಖನಿ} mean? (\textquotedblleft{}A pen\textquotedblright{}.) Say it once more.
\end{tcolorbox}
\section[\texorpdfstring{\kn{ಕಾಗದ} (kāgada)}{ಕಾಗದ (kāgada)}]{\kn{ಕಾಗದ} (kāgada) — paper; a letter}

\label{lesson:KA-C118-kagada}



\begin{quote}
Before the new one: say the Kannada for a bag, then the Kannada for a pen.

Read the lines aloud as one run: \textbf{\kn{ಇದು} \kn{ಬಾಗಿಲು}. \kn{ಕುರ್ಚಿ} \kn{ಇಲ್ಲಿ} \kn{ಇದೆ}.}
\end{quote}

\subsection*{You'll want to know: \kn{ಕಾಗದ}}
\textbf{\kn{ಕಾಗದ}} — \emph{kāgada} — \textquotedblleft{}paper; a letter\textquotedblright{}.

\begin{grammarlens}[title={What you've built}]
One more word for this chapter.
\end{grammarlens}

\subsection*{Your turn}
\begin{itemize}
\raggedright
  \item \emph{Say it:} \emph{kāgada}
  \item \emph{Say it:} \emph{kāgada}, once more
  \item \emph{Say it:} say \emph{kāgada}
  \item \emph{Recall:} say the Kannada for a bag, then the Kannada for a pen, then say \emph{kāgada} again
\end{itemize}

\begin{tcolorbox}[breakable,colback=teal!4,colframe=teal!35!black,title={Before you move on}]
What does \kn{ಕಾಗದ} mean? (\textquotedblleft{}Paper; a letter\textquotedblright{}.) Say it once more.
\end{tcolorbox}
